For a (not necessarily reduced) plabic graph $G$ we let $\Lplab_G$ be the corresponding plabic graph link, defined as in \cref{sec:intro:positroid_links}. Following~\cite{GL_qtcat}, we let  $\Ptop_L(q)$ be obtained from the top $a$-degree term of $\HOMP(L;a,z)$ by substituting $a:=q^{-\frac12}$ and $z:=q^{\frac12}-q^{-\frac12}$. The following conjecture generalizes~\cite[Theorem~1.11]{GL_qtcat} (see also~\cite{STZ,STWZ}).
\begin{conjecture} \label{conj:main}
Let $G$ be a simple plabic graph with $\conn(G)$ connected components.  Then we have
\begin{equation}\label{eq:FLY=pcnt}
  \RQG(q)= (q-1)^{\conn(G)-1}\Ptop_{\Lplab_G}(q).
\end{equation}
%
\end{conjecture}
For examples, see \cref{fig:intro_ex}, \cref{ex:Dynkin}, and \cref{sec:conjectures-examples}.

When the plabic graph $G$ is reduced, \cref{conj:main} becomes~\cite[Theorem~1.11]{GL_qtcat}. We present a different proof in \cref{sec:sub_skein-relation} by giving a plabic graph interpretation of the \FLY skein relation~\eqref{eq:HOMFLY_dfn}. More generally, in \cref{sec:leaf-recurr-plab}, we introduce a class of \emph{leaf recurrent plabic graphs} and prove \cref{conj:main} for them. 

%
 %
%
%
%
\begin{theorem}\label{thm:FLY=pcnt_leaf_rec}
  \Cref{conj:main} holds for leaf recurrent plabic graphs.
\end{theorem}
 %
 The fact that reduced plabic graphs are leaf recurrent was shown in~\cite[Remark~4.7]{MSLA}. Another natural subclass of leaf recurrent plabic graphs are the \emph{plabic fences} of~\cite[Section~12]{FPST}.
%
 A \emph{plabic fence} is a plabic graph obtained by drawing $k$ horizontal strands and inserting an arbitrary number of black-white and white-black bridges between them; see \cref{fig:plabic_fence} for an example and \cref{sec:plabic-fences} for a precise definition.
%
\begin{figure}
  \includegraphics[width=0.5\textwidth]{figures/plabic_fence}
  \caption{\label{fig:plabic_fence} A plabic fence.}
\end{figure}


\begin{proposition}\label{prop:reduced_and_fence_are_leaf_rec}
Reduced plabic graphs and plabic fences are leaf recurrent. In particular, \cref{conj:main} holds for these classes of plabic graphs.
\end{proposition}
%
This result gives a new proof of~\cite[Theorem~1.11]{GL_qtcat}.

A primary difficulty in showing \cref{conj:main} in full generality comes from the examples described in \cref{sec:notleaf,sec:notLA}. \Cref{sec:notleaf} explains that the simple plabic graph $G$ in \figref{fig:simple}(a) is not leaf recurrent. The quiver $\QG$ is \Louise (in fact, it is mutation equivalent to an acyclic quiver), and thus the point count polynomial $\RQG(q)$ can be easily checked to coincide with $\Ptop_{\Lplab_G}(q)$. \cref{sec:notLA} describes a simple plabic graph $G$ such that $\QG$ is not locally acyclic, in which case the task of computing $\RQG(q)$ becomes much harder.
%

Both sides of~\eqref{eq:FLY=pcnt} are rational functions in $q$ which are specializations of rational functions in $(q,t)$ arising from the cohomology of cluster varieties on the one hand and Khovanov--Rozansky link homology on the other hand. Thus, one can make a stronger conjecture where both sides of~\eqref{eq:FLY=pcnt} become polynomials in $(q,t)$. We discuss this in more detail in \cref{sec:clust-cohom-vs-link-hom}.
% We discuss these more general conjectures in \cref{sec:clust-cohom-vs-link-hom}.
%

\begin{remark}
In special cases, the cluster point count rational function $R(Q;q)$ has an explicit combinatorial description as a positive sum of terms of the form $(q-1)^a q^b$.  This holds, for example, when $Q$ is acyclic \cite{LS2} (see \cref{prop:acycliccount}), or when the cluster variety is a Richardson variety \cite{Deodhar}, see also \cite[Section 9]{GL_qtcat}.    We remark that such formulae should be compared to ruling polynomials of links \cite{Ruth}, though we do not know such a formula for general plabic graph quivers $Q_G$.
\end{remark}

\section{Positroid combinatorics}\label{sec:positr-comb}
In this section, we review some background on the combinatorics of plabic graphs and the associated objects.


\subsection{Bounded affine permutations}\label{sec:bound-affine-perm}
A \emph{$(k,\n)$-bounded affine permutation}~\cite{KLS} is a bijection $f:\Z\to\Z$ such that 
\begin{itemize}
\item $f$ is \emph{$\n$-periodic}: $f(i+\n)=f(i)+\n$ for all $i\in \Z$,
\item  $i\leq f(i)\leq i+\n$ for all $i\in\Z$, and %
\item $\sum_{i=1}^\n(f(i)-i)=k\n$.
\end{itemize}
The set of $(k,\n)$-bounded affine permutations is denoted $\Boundkn$. Each $f\in\Boundkn$ gives rise to a unique permutation $\fbar\in\Sn$ defined by $\fbar(i)\equiv f(i)\pmod \n$ for each $i\in[\n]$. Given $f\in\Boundkn$, an integer $i\in\Z$ is called a \emph{loop} (resp. a \emph{coloop}) of $f$ if $f(i)=i$ (resp. $f(i)=i+\n$). Thus, $\fbar$ has no fixed points if and only if $f$ has no loops and no coloops. In this case, $f$ is uniquely determined by $\fbar$, and the integer $k$ is recovered from $\fbar$ as follows:
\begin{equation*}%
  k=k(\fbar):=\#\{i\in[\n]\mid \fbar(i)<i\}.
\end{equation*}
We define $\taukn\in\Boundkn$  by
\begin{equation*}%
  \taukn(i)=
  \begin{cases}
    i, &\text{if $1\leq i\leq \n-k$,}\\
    i+\n, &\text{if $\n-k+1\leq i\leq \n$,}
  \end{cases}
\end{equation*}
and extend it to a function $\taukn:\Z\to\Z$ by $\n$-periodicity. 


Let $\Qkn:=\{(v,w)\in \Sn\times \Skn\mid v\leq w\}$, where $\leq$ denotes the Bruhat order on $\Sn$. Given a permutation $u\in\Sn$, we can extend it to an $\n$-periodic map $\tilde u:\Z\to\Z$ defined by $\tilde u(i)=u(i)$ for $i\in[\n]$. For $(v,w)\in\Qkn$, let $f_{v,w}:=\tilde w\taukn\tilde v^{-1}$. By~\cite[Proposition~3.15]{KLS}, the map $(v,w)\mapsto f_{v,w}$ gives a bijection $\Qkn\xrightarrow{\sim} \Boundkn$. Taking the inverse of this map, we obtain the following result, justifying the definition of a Richardson link in \cref{sec:intro:Rich_link}.
\begin{corollary}\label{cor:KLS}
If $\fp\in\Sn$ has no fixed points then it can be uniquely factored as a product $\fp=wv^{-1}$, where $(v,w)\in\Qkn$ and $k=k(\fp)$.
\end{corollary}

The open positroid stratification~\cite{KLS} of the Grassmannian contains a unique open dense stratum. It is labeled by the bounded affine permutation $\fkn\in\Boundkn$ sending $i\mapsto i+k$ for all $i\in\Z$. We let $\pikn\in\Sn$ be the permutation sending $i\mapsto i+k$ modulo $\n$ for all $i\in[\n]$.


\subsection{Plabic graphs}\label{sec:plabic-graphs}
%
%
%
%
%
%
%
%
%
%
%

\begin{figure}
  \includegraphics[width=0.7\textwidth]{figures/moves}
  \caption{\label{fig:moves} Local moves for plabic graphs: (a) contraction-uncontraction move; (b) square move; (c) tail addition/removal.}
\end{figure}




%

%

Let $G$ be a plabic graph. Recall from \cref{sec:intro:plabic_link} that $G$ must be nonempty, that the boundary vertices of $G$ are labeled by $1,2,\dots,\n$ in clockwise order, and that the strand permutation $\pig:[\n]\to[\n]$ sends $i\mapsto j$ whenever the strand starting at $i$ terminates at $j$. We always assume that our plabic graphs $G$ are such that $\pig$ has no fixed points, although most of our constructions can be easily extended to this case; see \cref{sec:fixed_pts}.

%

\begin{definition}[\cite{Pos}]
We say that $G$ is \emph{reduced} if it has the minimal number of faces among all plabic graphs with strand permutation $\pig$.
\end{definition}
Alternatively~\cite[Theorem~13.2]{Pos}, a plabic graph $G$ is reduced if and only if $G$ has no closed strands, no self-intersecting strands, and no pairs of strands forming a \emph{bad double crossing} shown in \figref{fig:double}(right). \emph{Good double crossings} shown in \figref{fig:double}(left) are allowed.

We consider several types of \emph{local moves} on plabic graphs, shown in \cref{fig:moves}. 
%
 It was shown in~\cite{Pos} that every permutation $\pi\in\Sn$ is the strand permutation of some reduced plabic graph $G$, and that moreover, any two reduced plabic graphs with the same strand permutation are related by a sequence of the moves (a) and (b) in \cref{fig:moves}. Below we review one construction of plabic graphs which will be particularly useful to us.

\subsection{Le-diagrams}\label{sec:Le-diags}
Let $\la=(\la_1\geq\la_2\geq\cdots\geq\la_k\geq0)$ be a partition, which we identify with its Young diagram. We assume that this Young diagram fits inside a $k\times (\n-k)$ rectangle, i.e.  satisfies $\la_1\leq \n-k$. We draw Young diagrams using \emph{English} (or \emph{matrix}) \emph{notation}. We use matrix coordinates for the boxes of $\la$, thus, row $i$ contains boxes with coordinates $(i,j)$ for $1\leq j\leq \la_i$, and box $(1,1)$ is the top left box.

A \emph{Le-diagram} $\Gamma$ \emph{of shape $\la$} is a way of placing a dot in some of the boxes of $\la$ so that every box that is below a dot in the same column and to the right of a dot in the same row must also contain a dot.  
%
 An example of a Le-diagram of shape $\la=(3,3,3,1)$ is shown in \figref{fig:Gamma-to-L'}(left). Le-diagrams whose shape fits inside a $k\times(\n-k)$ rectangle are in bijection with the elements of $\Qkn$ and $\Boundkn$; see~\cite[Section~19]{Pos}. Given a Le-diagram $\Gamma$ of shape $\la$, one can read off the corresponding pair $(v,w)\in\Qkn$ as follows. First, $w\in\Skn$ is the $k$-Grassmannian permutation corresponding to $\la$. Specifically, let us label the unit steps in the southeast boundary of $\la$ by $1,2,\dots,\n$ in increasing order from the northeast to the southwest corner. Then the labels of the vertical steps form a $k$-element subset of $[\n]$, and this set is precisely $\{w(\n-k+1),\dots,w(\n)\}$. This determines $w$ uniquely. Alternatively, we can label each box $(i,j)$ of $\la$ with the simple transposition $s_{k+j-i}$, and then a reduced word for $w$ is obtained by reading these simple transpositions in the northwest direction. Thus, any reduced word for $w$ ends with $s_k$ which labels the box $(1,1)$. To obtain a reduced word for $v$, we read these simple transpositions but ignore the ones labeled by dots in $\Gamma$. Conversely, given $(v,w)\in\Qkn$, the shape $\la$ of $\Gamma$ is reconstructed from $w$, and the empty boxes of $\Gamma$ correspond to the rightmost subexpression for $v$ inside the (unique up to commutation) reduced word for $w$.
\begin{example}
For the Le-diagram in \figref{fig:Gamma-to-L'}(left), we have 
\begin{equation*}%
w= s_7 s_3 s_4 s_5 s_6 s_2 s_3 s_4 s_5 s_1 s_2 s_3 s_4
\quad\text{and}\quad
v= s_6 s_3 s_2.
\end{equation*}
\end{example}

\begin{figure}
  \includegraphics[width=0.7\textwidth]{figures/Le_to_G}
  \caption{\label{fig:Le_to_G} Converting a Le-diagram (top row) into a plabic graph (middle row) and into a link (bottom row).}
\end{figure}



A Le-diagram $\Gamma$ can be converted into a plabic graph $G(\Gamma)$ using the local rules shown in the first two rows of \cref{fig:Le_to_G}. This plabic graph is always reduced and has strand permutation $\fp=wv^{-1}$, where $(v,w)\in\Qkn$ are recovered from $\Gamma$ using the above procedure.


%

\subsection{Properties of plabic graph links}\label{sec:plabic_links_properties}
Let $G$ be a reduced plabic graph with strand permutation $\pi=\pig$. Recall the description of the plabic graph link $\Lplab_\pi=\Lplab_G$ from \cref{sec:intro:plabic_link}. It may appear from this description that $\Lplab_G$ depends on the precise way of drawing the strands in $G$ as smooth curves, or that $\Lplab_G$ changes if one rotates the graph $G$ (since the over/under-crossings information depends on the complex arguments of the strands). However, it turns out that the link $\Lplab_G$ is in fact independent of these choices, as explained in~\cite{ACampo,STWZ,FPST}. Let us give a more invariant description of $\Lplab_G$.

Recall that $G$ is embedded in a disk $\disk:=\{x\in\C: |x|\leq 1\}$. Consider the solid torus $\disk\times S^1$, and consider an equivalence relation $\sim$ on it defined as follows: for $(x,y),(x',y')\in \disk\times S^1$, write $(x,y)\sim(x',y')$ if $x,x'\in\partial \disk$ belong to the boundary of the disk and $x=x'$. In other words, we collapse each fiber $x\times S^1$, where $x\in\partial \disk$, to a point. The resulting quotient space $\disk\times S^1/\sim$ is homeomorphic to a $3$-dimensional sphere $S^3$.

An \emph{oriented divide} is a smooth immersion of a union of oriented circles (called \emph{branches}) into $\disk$. 
 %
%
%
%
%
 Divides are required to satisfy certain genericity conditions, e.g. that any intersections between the branches must be transversal and belong to the interior of $\disk$, and that there have to be no triple intersections; see~\cite[Definitions~2.1 and~8.1]{FPST} for a complete list. 

Given a branch $\gamma$ of an oriented divide $\Divide$, we lift each point of $\gamma$ to $\disk\times S^1$ by letting the $S^1$ coordinate be the complex argument of the tangent vector of $\gamma$ at this point. By the genericity assumptions, we obtain an embedding of a union of oriented circles into $S^3$, i.e.  a link, denoted $L_\Divide$.

\begin{figure}
  \includegraphics[width=0.7\textwidth]{figures/divide_moves}
  \caption{\label{fig:divide_moves} Local moves for divides~\cite[Figures~28--30]{FPST}.}
\end{figure}


The link $L_\Divide$ is invariant under applying the local moves shown in \cref{fig:divide_moves} to $\Divide$; see~\cite[Proposition~8.4]{FPST}. We thank the anonymous referee for pointing out that the map $\Divide \mapsto L_\Divide$ coincides with the map reconstructing a Legendrian curve from its front projection. Thus, the invariance under local moves may be deduced from e.g.~\cite{Arnold}.

Take the strands in our reduced plabic graph $G$. Each of the $\n$ boundary vertices of $G$ has one incoming strand and one outgoing strand. Identifying their endpoints, we obtain an oriented divide $\Divide(G)$. It was shown in~\cite{Hirasawa} that the associated link $L_{\Divide(G)}$ is isotopic to the link $\Lplab_G$ described in \cref{sec:intro:plabic_link}. This shows that the link $\Lplab_G$ is indeed invariant under rotation of $G$ and under small perturbations of the branches of $\Divide(G)$, since this is clearly the case for $L_{\Divide(G)}$. 

\begin{remark}\label{rmk:rotation}
In what follows, it will be convenient to us to use rotational invariance of $\Lplab_G$. In our figures, we will fix some angle $\alpha\in[0,2\pi)$, and assume that the link $\Lplab_G$ is obtained by first rotating $G$ by the angle $-\alpha$, then applying the description from \cref{sec:intro:plabic_link}, and then rotating the picture back by $\alpha$. In other words, if two strands $S_1,S_2$ in $G$ intersect at some point $p$, we consider their complex arguments of tangent vectors in a different interval: $\Arg(S_1,p),\Arg(S_2,p)\in[\alpha,\alpha+2\pi)$. Then we draw $S_1$ above $S_2$ if and only if $\Arg(S_2,p)>\Arg(S_1,p)$. And then for all points $p$ where some strand $S$ satisfies $\Arg(S,p)=\alpha$, we insert line segments going to the boundary and back. In the figures, the direction $\alpha$ is indicated by the ``over/under circle'' \\
%
\centerline{\includegraphics[width=0.17\textwidth]{figures/over_under_2}}
%
and each point $p$ on a strand $S$ satisfying $\Arg(S,p)=\alpha$ is marked by $\mysq$ as in \cref{fig:over_under}.
\end{remark}






\section{Positroid link isotopies}\label{sec:positr-link-isot}
The goal of this section is to prove \cref{thm:isotopy}. We first discuss some properties of plabic graph links, and then we construct a sequence of isotopies, for any permutation $\pi\in\Sn$ without fixed points, of the links
\begin{equation*}%
\Lrich_\pi \cong  \Lplab_\pi \cong \Lperm_\pi \cong \Ltor_\pi,
\end{equation*}
in that order. When $\pi=\pi_C$ for some concave curve $C$, we will also show that
\begin{equation*}%
  \Ltor_\pi\cong\Lcox_\pi,
\end{equation*}
completing the proof.

%
%

Throughout this section, we fix a permutation $\pi\in\Sn$ without fixed points, and let $(v,w)\in\Qkn$ be the corresponding pair satisfying $\pi=wv^{-1}$,  $\Gamma$ the corresponding Le-diagram of shape $\la$, and $G$ the corresponding plabic graph, obtained from $\Gamma$ using the recipe in the first two rows of \cref{fig:Le_to_G}.

\subsection{From Richardson links to plabic graph links}\label{sec:from-rich-links}
 Our goal is to construct an isotopy between the Richardson link $\Lrich_{v,w}$ (\figref{fig:Lrich}(right)) and the plabic graph link $\Lplab_G$ (\figref{fig:Lplab}(c)).

Recall from \cref{rmk:Rich_Young} that we are drawing the Richardson braid $\beta_{v,w}=\beta(w)\cdot \beta(v)^{-1}$ in a particular way, so that the crossings of $\beta(w)$ form a shape obtained by rotating $\la$ clockwise by $135^\circ$, and the crossings of $\beta(v)^{-1}$ are drawn in some boxes of a shape obtained from the $135^\circ$-rotation of $\la$ by reflecting it along the vertical axis; see \figref{fig:Lrich}(left). Moreover, the crossings of $\beta(v)^{-1}$ are located precisely in the boxes of $\la$ which do not contain a dot in $\Gamma$. 
 %

\begin{figure}
  \includegraphics[width=1.0\textwidth]{figures/Lrich_to_Lprime}
  \caption{\label{fig:Lrich-to-L'} Converting the link $\Lrich_{v,w}$ into $\Lrichp_{v,w}'$; see \cref{sec:from-rich-links}.}
\end{figure}


Our first goal is to draw the link $\Lrich_{v,w}$ entirely within $\la$. To do that, we take $\beta(v)^{-1}$, reflect it along the vertical axis (flipping the over/under-crossings), and place the resulting braid on top of $\beta(w)$. We rotate the result counterclockwise by $135^\circ$. Every line segment in the boundary of $\la$ contains an endpoint of a strand in $\beta(v)^{-1}$ and an endpoint of a strand in $\beta(w)$. We join these strand endpoints, obtaining a link diagram drawn inside of $\la$. This link $\Lrichp_{v,w}'$, shown in \cref{fig:Lrich-to-L'}, is clearly isotopic to $\Lrichp_{v,w}$. Intuitively, the transformation $\Lrich_{v,w}\to\Lrichp_{v,w}'$ can be described as follows: open a book, and place the left (resp. right) half of \figref{fig:Lrich}(right) onto the left (resp. right) page of the book. Then close the book with your right hand, so that the front cover of the book is at the bottom. The page containing (the reflection of) $\beta(v)^{-1}$ will be on top of the page containing $\beta(w)$. Finally, rotate the closed book $135^\circ$ counterclockwise. 


\begin{figure}
  \includegraphics[width=0.9\textwidth]{figures/Gamma_to_Lprime}
  \caption{\label{fig:Gamma-to-L'} Converting a Le-diagram $\Gamma$ into a link $L'_{v,w}$.}
\end{figure}




An alternative description of the link $\Lrichp_{v,w}'$ can be obtained directly from the Le-diagram $\Gamma$ by replacing each box containing a dot as shown in \figref{fig:Gamma-to-L'}(middle top) and each box not containing a dot as shown in \figref{fig:Gamma-to-L'}(middle bottom). For example, the Le-diagram in \figref{fig:Gamma-to-L'}(left) gets converted into the link $\Lrichp_{v,w}'$ in \figref{fig:Gamma-to-L'}(right); this is the same link as in \figref{fig:Lrich-to-L'}(right). Since $\pi$ has no fixed points, each row and each column of $\Gamma$ contains at least one dot. Consider some column $j$ of $\Gamma$ and let $b_j$ be the highest box in it which contains a dot. The part of $\Lrichp_{v,w}'$ above $b_j$ contains a strand going vertically to the top boundary and then back, passing between all the horizontal strands that it encounters along the way. We may therefore contract this piece of the strand to be fully contained inside $b_j$, as shown in \figref{fig:pull}(far left). Repeating this procedure for each column $j=1,2,\dots,\la_1$, we obtain a link $\Lrichp_{v,w}''$ isotopic to $\Lrichp_{v,w}'$; see \figref{fig:pull}(middle).

\begin{figure}
  \setlength{\tabcolsep}{0pt}
\begin{tabular}{c|c}
\begin{tikzpicture}[baseline=(Z.base)]
\coordinate(Z) at (0,0);
\node(A) at (0,2.05){
  \includegraphics[width=0.2\textwidth]{figures/pull}
};
\end{tikzpicture}
&
  \includegraphics[width=0.78\textwidth]{figures/pull_to_L}
\end{tabular}
  \caption{\label{fig:pull} An isotopy from $\Lrichp_{v,w}'$ to $\Lplab_G$.}
\end{figure}



Finally, we claim that the link $\Lrichp_{v,w}''$ is isotopic to the plabic graph link $\Lplab_G$. To see this, we choose the angle $\alpha$ from \cref{rmk:rotation} to be $80^\circ$. Then combining the rules from \cref{fig:Le_to_G} for converting the Le-diagram $\Gamma$ into the plabic graph $G$ with the description of $\Lplab_G$ given in \cref{sec:plabic_links_properties}, we see that a planar diagram of $\Lplab_G$ is obtained from $\Gamma$ via the rules shown in the last two rows of \cref{fig:Le_to_G}. Replacing each square (where $\Arg(S,p)=\alpha$) by a horizontal segment going to the left boundary above all other strands and back below all other strands, we see that $\Lplab_G= \Lrichp_{v,w}''$; see \figref{fig:pull}(far right). To summarize, we have constructed explicit isotopies
\begin{equation*}%
  \Lrich_\pi=\Lrich_{v,w}\cong\Lrichp_{v,w}'\cong\Lrichp_{v,w}''=\Lplab_G.
\end{equation*}



\begin{figure}
  \includegraphics[width=0.5\textwidth]{figures/double}
  \caption{\label{fig:double} Good and bad double crossings of strands in plabic graphs.}
\end{figure}


\subsection{From plabic graph links to permutation links}
Our goal is to find an isotopy $\Lplab_\pi\cong \Lperm_\pi$. Let $G$ be a reduced plabic graph with strand permutation $\pi$. Let $\Divide=\Divide(G)$ be the oriented divide obtained from $G$ as in \cref{sec:plabic_links_properties}. Let $1,2,\dots,\n$ be the boundary vertices of $G$. Since $G$ is reduced, it contains no bad double crossings shown in \figref{fig:double}(right). Thus, it is clear that applying the moves from \cref{fig:divide_moves}, one can transform $\Divide$ into an oriented divide $\Divide'$ obtained by drawing a straight arrow $i\to \pi(i)$ for each $i\in[\n]$. The moves from \cref{fig:divide_moves} give rise to link isotopies, so $L_\Divide\cong L_{\Divide'}$. Let us move the boundary vertices $1,2,\dots,\n$ smoothly to the points $1',2',\dots,\n'$ which are located clockwise on the left semicircle of $\partial \disk$, i.e.  on the subset of $\partial\disk$ with negative real part. Let $\Divide''$ be the oriented divide obtained by drawing a straight arrow $i'\to \pi(i)'$ for each $i\in[\n]$. It follows that $L_{\Divide'}\cong L_{\Divide''}$.
%
 Let $i,j\in[\n]$ be such that the corresponding arrows $S_i:=(i'\to \pi(i)')$ and $S_j:=(j'\to \pi(j)')$ of $\Divide''$ cross at some point $p$. Consider the arguments of the tangent vectors $0\leq \Arg(S_i,p)\neq\Arg(S_j,p)<2\pi$. Then it is easy to check that $\Arg(S_i,p)>\Arg(S_j,p)$ if and only if $\pi(i)<\pi(j)$. 
%
 Comparing to the description in \cref{sec:intro:perm_links}, we get that $L_{\Divide''}\cong\Lperm_\pi$.



\subsection{From permutation links to toric permutation links}\label{sec:perm-to-toric}
Our goal is to find an isotopy $\Lperm_\pi\cong\Ltor_\pi$. Recall from \cref{rmk:grid} that $\Ltor_\pi$ is isotopic to the \emph{planar grid link} denoted $\Lgrid_\pi$; see \figref{fig:grid_iso}(a--c). In $\Lgridt_\pi$ and $\Lgrid_\pi$, the vertical arrows are drawn above the horizontal arrows. 

The main diagonal $y=1-x$ of the square $[0,1]^2$ divides $\Lgrid_\pi$ into the \emph{lower-triangular} and the \emph{upper-triangular} parts. In the lower-triangular part, all vertical arrows point down and all horizontal arrows point right, while in the upper-triangular part, all vertical arrows point up and all horizontal arrows point left. 
Let us take the lower-triangular part and reflect it around the main diagonal, placing it below the upper-triangular part; see \figref{fig:grid_iso}(c--d). We obtain a planar diagram of a link which we denote $L_\pi'$. Clearly, $L_\pi'\cong \Lgrid_\pi$. The diagram of $L_\pi'$ contains arrows pointing, up, left, down, and right. It follows that the over/under-crossings rule  in $L_\pi'$ is given by the following total order on the arrow directions: up $>$ left $>$ down $>$ right. In particular, $L_\pi'$ is a divide link for the angle $\alpha=45^\circ$ from \cref{rmk:rotation} (where the main diagonal is treated as part of the boundary of the disk), and we get that $L_\pi'\cong\Lperm_\pi$. 
%

\subsection{From toric permutation links to Coxeter links}\label{sec:tor_to_Cox}
First, we introduce a more general class of links, associated to arbitrary monotone curves inside the $k\times(\n-k)$ rectangle.

%
%
We say that $C_g$ is a \emph{monotone curve} if it is the plot of a strictly monotone increasing function $g(x):[0,\n-k]\to[0,k]$ satisfying $g(0)=0$, $g(\n-k)=k$, and $g(j)\notin\Z$ for all $j=1,2,\dots,\n-k-1$. Thus, the curve $C_g$ passes through the points $(0,0)$ and $(\n-k,k)$, but through no other lattice points. To any monotone curve $C_g$ we associate a link $L_g$ drawn on the surface of $\T$. The planar diagram of $L_g$ is obtained from the projection $\Cproj_g$ of $C_g$ to $\T$ via the following rule: whenever  two points $(x,g(x))$, $(x',g(x'))$ on $C_g$ (with $x<x'$) project to the same point of $\Cproj_g$, we draw the projection of $(x,g(x))$ above the projection of $(x',g(x'))$. 

The link diagram of $L_g$ is drawn on the torus $\T$, and therefore it is natural to think of $L_g$ itself as a link inside the \emph{thickened torus} $\T\times[0,1]$. Given a point $p=(x,g(x))\in C_g$, we denote by $\projtx p\in \T$ the projection of $p$ to $\T$, and we let $\left(\projtx p,1-\frac{x}{\n-k}\right)$ be the corresponding point in $\T\times[0,1]$. This yields an embedding of $C_g$ into $\T\times[0,1]$. Under this embedding, the points $(0,0)$ and $(\n-k,k)$ map to $((0,0),1)$ and $((0,0),0)$, respectively. Adding a line segment connecting $((0,0),1)$ to $((0,0),0)$, we obtain a representation of the link $L_g$ inside $\T\times[0,1]$. An obvious consequence of this construction is that the link $L_g$ only depends on the set of lattice points below the plot of $g$.
\begin{corollary} \label{cor:monotone_curves}
Let $C_g$, $C_h$ be two monotone curves, and assume that for each $j=1,2,\dots,\n-k-1$, we have $\lfloor g(j)\rfloor=\lfloor h(j)\rfloor$. Then the links $L_g$ and $L_h$ are isotopic. 
\end{corollary}

Suppose now that $\pi=\pi_C$ for a concave curve $C$. Our goal is to find an isotopy $\Ltor_\pi\cong\Lcox_C$. Comparing the above description to the one in \cref{sec:intro:cox_link}, we see that when $g$ is the concave down function satisfying $C=C_g$, we have $\Ltor_\pi\cong L_g$. Our goal is to choose a monotone curve $C_h$ such that $L_g\cong L_h$ by \cref{cor:monotone_curves}, and such that $L_h\cong \Lcox_C$. 
 %

\begin{figure}
\includegraphics[width=1.0\textwidth]{figures/C_to_beta}
  \caption{\label{fig:C_to_beta} Converting a concave curve $C$ into the Coxeter link $\beta(\abf(C))$; see \cref{sec:tor_to_Cox}.}
\end{figure}


Let $\la=\la_C=(\la_1\geq \la_2\geq\dots\geq \la_k=0)$ be the Young diagram above $C$. For $r=0,1,\dots,k-1$, let $p_r:= \left(\la_{k-r}+\frac rk,r\right)$. Thus, $p_0=(0,0)$. Let $p_{k}:=(\n-k,k)$. Let $h$ be the piecewise-linear function whose plot passes through the points $p_0,p_1,\dots,p_k$, and let $C_h$ be the corresponding monotone curve; see \figref{fig:C_to_beta}(b). Indeed, by \cref{cor:monotone_curves}, we have $L_g\cong L_h$. It remains to establish the isotopy $L_h\cong \Lcox_C$. 

Let $\beta=\beta(\abf(C))$ be the Coxeter braid associated to $C$, defined in~\eqref{eq:beta_abf_dfn}. It is a braid on $k$ strands. For $i=1,2,\dots,k$, let $S_i(\beta)$ be the strand of $\beta$ whose left endpoint is labeled by~$i$. Thus, the right endpoint of $S_i(\beta)$ is labeled by $i+1$ for $i<k$ and by $1$ for $i=k$. On the other hand, the link diagram of $L_h$ in $\T$ intersects the line $y=0$ in exactly $k$ points with horizontal coordinates $\frac rk$ for $r=0,1,\dots, k-1$. We would like to view $L_h$ as the closure of a $k$-strand braid $\beta_h$ in the solid torus; see \figref{fig:C_to_beta}(c--d). The braid $\beta_h$ is obtained from $L_h$ via the following procedure: whenever a strand of $L_h$ passes through some point $(1,y)\in[0,1]^2$ and continues from the point $(0,y)\in[0,1]^2$, we connect the points $(0,y)$ and $(1,y)$ by a horizontal line segment that is drawn below all strands of $L_h$. The result is a drawing of a braid connecting $k$ points at the bottom of $[0,1]^2$ to $k$ points at the top of $[0,1]^2$. 

For $i=1,2,\dots, k-1$, let $S_i(\beta_h)$ be the piece of $\beta_h$ connecting the point $(\frac{i-1}k,0)$ to the point $(\frac ik,1)$. For $i=k$, let $S_k(\beta_h)$ be the remaining piece of $\beta_h$, connecting $(\frac{k-1}k,0)$ to $(0,1)$. See \figref{fig:C_to_beta}(d--e).

We claim that the braids $\beta_h$ and $\beta$ are isotopic. We compare them strand-by-strand, identifying $S_i(\beta)$ with $S_{k+1-i}(\beta_h)$ for $i=1,2,\dots,k$, in that order. (This correspondence is indicated via colors in \figref{fig:C_to_beta}(d--e).) For convenience, we rotate the drawing of $\beta$ by $90^\circ$ counterclockwise. The strand $S_1(\beta)$ moves straight up during the $\ell_2^{a_2}\cdots \ell_k^{a_k}$ part, and then it moves left under all other strands. We apply an isotopy to make $S_k(\beta_h)$ move straight up towards the top boundary of $[0,1]^2$, and then move left under all other strands. Recall that we have a sequence $\abf(C)=(a_2,\dots,a_k)$ given by $a_i=\la_{i-1}-\la_i$ for $i=2,\dots,k$. Thus, the strand $S_2(\beta)$ wraps around $S_1(\beta)$ exactly $a_2$ times. On the other hand, $S_{k-1}(\beta_h)$ is the projection to $\T$ of the line segment of $C_h$ connecting $p_{k-1}$ to $p_k$. This projection intersects the vertical boundary of $\T$ exactly $a_2$ times. Therefore $S_{k-1}(\beta_h)$ wraps around $S_k(\beta_h)$ exactly $a_2$ times. Continuing in this fashion, we see that for $i=3,\dots,k$, the strand $S_i(\beta)$ wraps around the union $S_1(\beta)\cup S_2(\beta)\cup\dots\cup S_{i-1}(\beta)$ exactly $a_i$ times. On the other hand, the strand $S_{k+1-i}(\beta_h)$ wraps around the union $S_k(\beta_h)\cup S_{k-1}(\beta)\cup\dots\cup S_{k+2-i}$ exactly $a_i$ times. This shows that the braids $\beta$ and $\beta_h$ are isotopic, and therefore we get an isotopy $\Lcox_C\cong L_h$ between their closures.

\begin{remark}
The above construction is extended to links with multiple components (resp. to monotone curves passing through some lattice points in the interior of the rectangle) in~\cite[Section~7.3]{GL_EHA}.
\end{remark} 


